% TOP_PROBLEM: {"id":30007605,"problem_number":"LOCAL-30007605","title":"Macroprudential leakage","category_id":21,"category":{"id":21,"name":"economics","display_name":"Economics","description":"Open economic theory statements, published research agendas, and proposed scoped specifications; question provenance and historical priority are qualified per record.","slug":"economics","order_index":21,"created_at":"2026-10-08T00:00:00Z"},"status":"research_question","external_url":null,"legacy_ids":[],"published":false,"statement":"In the explicitly specified setting: How much do bank-focused restrictions shift risky credit to nonbanks, other sectors or foreign lenders? The mathematical statement above fixes the target, assumptions and scope of this catalogue formulation.","economics":{"original_id":94,"field":"Banking, credit and financial stability","kind":"identification","formalization_basis":"adapted_research_question","source_status":"research_agenda","priority_status":"earliest_found","priority_note":"This is the earliest explicit statement verified in this audit; first priority for the full catalogue formulation is not established.","earliest_verified_reference":{"authors":["Nina Biljanovska","Sophia Chen","Gaston Gelos","Deniz Igan","Maria Soledad Martinez Peria","Erlend Nier","Fabián Valencia"],"title":"Macroprudential Policy Effects: Evidence and Open Questions","year":2023,"url":"https://www.imf.org/en/publications/departmental-papers-policy-papers/issues/2023/03/26/macroprudential-policy-effects-evidence-and-open-questions-527293","doi":"10.5089/9798400226304.087","locator":"Overview; Core Questions; Key Findings and Themes, official publication landing-page summary","evidence_summary":"Asks leakage, resilience, activity and heterogeneous policy effects."},"current_reference":{"authors":["Nina Biljanovska","Sophia Chen","Gaston Gelos","Deniz Igan","Maria Soledad Martinez Peria","Erlend Nier","Fabián Valencia"],"title":"Macroprudential Policy Effects: Evidence and Open Questions","year":2023,"url":"https://www.imf.org/en/publications/departmental-papers-policy-papers/issues/2023/03/26/macroprudential-policy-effects-evidence-and-open-questions-527293","doi":"10.5089/9798400226304.087","locator":"Overview; Core Questions; Key Findings and Themes, official publication landing-page summary","evidence_summary":"Asks leakage, resilience, activity and heterogeneous policy effects."},"formalization_note":"Adapts an explicitly verified research-agenda passage. Mathematical objects, interventions, objectives and scope are proposed here; existing model-specific solutions are not relabeled unsolved."},"statusQualification":"Published question or research agenda verified at the cited locator. The mathematical specification is adapted for this catalogue and includes additional assumptions; source publication does not establish first-ever priority or certify this exact model remains unsolved. Adapts an explicitly verified research-agenda passage. Mathematical objects, interventions, objectives and scope are proposed here; existing model-specific solutions are not relabeled unsolved.","statusReviewedAt":"2026-10-08"}
% FIRST_POSTED: 2026-10-08
% ENTRY_KIND: statement_only
% !TeX program = lualatex
\documentclass[11pt]{article}
\usepackage{fontspec}
\usepackage[a4paper,margin=25mm]{geometry}
\usepackage{amsmath,amssymb,mathtools}
\usepackage{xurl}
\usepackage[hidelinks]{hyperref}
\newcommand{\sref}[2]{\hyperlink{source:#1:#2}{[S#2]}}
\setlength{\emergencystretch}{3em}
\begin{document}
% problemId:30007605
\section*{Macroprudential leakage}\label{30007605}
\subsection{Definitions and mathematical statement}
Fix a jurisdiction adopting a bank-focused macroprudential tightening \(u\), measured on a declared instrument scale. Let \(B_h(u)\) be regulated-bank risky credit at horizon \(h\geq0\), \(N_h(u)\) domestic nonbank credit, \(F_h(u)\) cross-border credit and \(S_h(u)\) credit redirected to unregulated sectors. Classifications must be mutually exclusive, or double-counted claims must be reconciled. For the intervention \(do(u)\), define the leakage ratio whenever its denominator is positive:
\[\Lambda_h=\frac{E[N_h(u)+F_h(u)+S_h(u)-N_h(0)-F_h(0)-S_h(0)]}{E[B_h(0)-B_h(u)]}.\]
The research task is to identify this ratio and the accompanying change in system-wide default or crisis risk. Specify the eligible borrowers, instrument, adoption date and untreated counterfactual; determine conditions under which policy endogeneity and anticipation can be controlled. Include lender and borrower adjustment, not only static accounting substitution. Estimate heterogeneity by instrument, financial development and stress state. If unregulated credit is only partially observed, derive bounds using transparent coverage assumptions. Credit displacement is not necessarily risky leakage: changes in borrower quality, collateral and maturity must be measured before interpreting \(\Lambda_h\) as a financial-stability effect.

\par\textbf{Formulation and scope.} Adapts an explicitly verified research-agenda passage. Mathematical objects, interventions, objectives and scope are proposed here; existing model-specific solutions are not relabeled unsolved.

\par\textbf{Open-question provenance.} An explicit question or research agenda was verified at the source locator. This does not by itself certify that every later version remains unresolved \sref{30007605}{1}.

\par\textbf{Historical priority.} This is the earliest explicit statement verified in this audit; first priority for the full catalogue formulation is not established.

\subsection{Short English statement}
In the explicitly specified setting: How much do bank-focused restrictions shift risky credit to nonbanks, other sectors or foreign lenders? The mathematical statement above fixes the target, assumptions and scope of this catalogue formulation.

\subsection{Sources}
\begin{itemize}\raggedright
\item[\textnormal{[S1]}]\hypertarget{source:30007605:1}{} Nina Biljanovska, Sophia Chen, Gaston Gelos, Deniz Igan, Maria Soledad Martinez Peria, Erlend Nier, Fabián Valencia. 2023. Macroprudential Policy Effects: Evidence and Open Questions. Overview; Core Questions; Key Findings and Themes, official publication landing-page summary.
\url{https://www.imf.org/en/publications/departmental-papers-policy-papers/issues/2023/03/26/macroprudential-policy-effects-evidence-and-open-questions-527293}.
DOI: 10.5089/9798400226304.087.
{\small\textit{Earliest verified question/agenda reference; historical priority qualified below. Asks leakage, resilience, activity and heterogeneous policy effects.}}
\end{itemize}
\end{document}
